\chapter{SONUÇLAR VE ÖNERİLER}
\label{ch:sonuclar}

Bu bölümde CloudMon sisteminin deneysel bulgularından çıkarılan sonuçlar özetlenmekte, çalışmanın özgün katkıları değerlendirilmekte ve çalışmanın sınırlılıkları ile gelecek araştırma önerileri tartışılmaktadır.


\section{Temel Bulgular}
\label{sec:sonuc-bulgular}

Bu çalışma; kapasite tahmini, açıklanabilirlik analizi ve log tabanlı erken uyarı olmak üzere üç ana araştırma ekseninde somut bulgular üretmiştir.

\subsection{Kapasite Tahmin Modeli}

42 model (7 algoritma $\times$ 3 kaynak $\times$ 2 ufuk) değerlendirmesinin en belirgin bulgusu, Doğrusal Regresyon'un (LR) ağaç tabanlı tüm algoritmaları geride bırakmasıdır.
24 saatlik ufukta LR, vCPU için R\textsuperscript{2}\,=\,0{,}9577, Bellek için R\textsuperscript{2}\,=\,0{,}7945 ve Disk için R\textsuperscript{2}\,=\,0{,}6984 değerlerine ulaşmıştır.
En yakın rakibi olan Gradyan Artırmalı Ağaçlar, vCPU 24\,s senaryosunda yalnızca R\textsuperscript{2}\,=\,0{,}1079 üretmiştir; bu fark, OpenStack zaman serilerinin test döneminde güçlü monoton trendler barındırdığını ve LR'nin bu trendi gecikme ve hareketli ortalama özellikleri üzerinden doğrudan öğrendiğini göstermektedir.

72 saatlik ufukta sonuçlar kaynak bazında ayrışmaktadır.
vCPU'da LR, R\textsuperscript{2}\,=\,0{,}7728 ile etkin tahmin gücünü korumuştur.
Bellek ve Disk'te ise MLP en iyi performer konumuna geçmiş; bununla birlikte R\textsuperscript{2} değerleri sırasıyla 0{,}3456 ve 0{,}2212'ye gerilemiştir.
Ağaç tabanlı modellerin büyük çoğunluğu 72 saatlik ufuklarda negatif R\textsuperscript{2} üretmiş, yani ortalama tahminden daha kötü performans göstermiştir.
Bu bulgu, söz konusu modellerin 272 özellikli yüksek boyutlu uzayda eğitim kümesine aşırı uyduğunu ve zaman serisi genellemesinde yetersiz kaldığını ortaya koymaktadır.

\subsection{SHAP Açıklanabilirlik Bulguları}

SHAP analizi, log tabanlı özelliklerin salt metrik özelliklerinden anlamlı biçimde farklılaştığını ve kapasite tahminine bağımsız katkı yaptığını nicel olarak göstermiştir.
vCPU 24\,s modelinin ilk 10 özelliği arasında dört log özelliği yer almaktadır:
\texttt{resource\_exhausted\_count\_rsum\_72h} (5.\ sıra, $|\phi|$\,=\,2\,150),
\texttt{total\_errors\_rsum\_72h} (6.\ sıra, $|\phi|$\,=\,1\,800),
\texttt{scheduler\_degradation\_count\_rsum\_72h} (9.\ sıra, $|\phi|$\,=\,1\,050) ve
\texttt{timeout\_count\_rsum\_72h} (10.\ sıra, $|\phi|$\,=\,1\,028).
Disk 24\,s modelinde ise \texttt{total\_errors\_rsum\_24h} birinci sıraya taşınmış; log özellikleri metrik özelliklerini bile geride bırakmıştır.
Bu bulgular, log verilerinin kapasite tahmininde salt gürültü değil, anlamlı bir sinyal kaynağı olduğunu kanıtlamakta ve hibrit özellik uzayının kurucu varsayımını doğrulamaktadır.

\subsection{Log Analizi ve Erken Uyarı}

Dört aşamalı log işleme pipeline'ı (trend tespiti, z-skor anomali tespiti, ön-arıza örüntü eşleştirme, cascade risk tespiti) test dönemi (Nisan--Haziran 2026) üzerinde anlamlı sinyaller üretmiştir.
\textit{Slow\_response} loglarının geçmiş 15 incident analizine göre arızadan ortalama 68 saat önce \%287 oranında artış gösterdiği bulunmuş ve mevcut test dönemi verisi bu örüntüyle örtüşmüştür.
Kota aşımı loglarının z-skoru 4{,}5'e çıkmış; disk I/O hatası (z\,=\,3{,}6), ağ hatası (z\,=\,2{,}8) ve veritabanı bağlantı hatası (z\,=\,2{,}6) da anormallik eşiğini (\,$\sigma > 2{,}5$) aşmıştır.
2026-06-28 tarihinde Nova, Neutron ve Keystone servislerinin eş zamanlı degradasyonu cascade risk kanalı tarafından tespit edilmiştir.
Bu bulgular, log tabanlı erken uyarı mekanizmasının metrik eşik alarmlarına kıyasla üç güne yakın önceden sinyal üretme potansiyeline sahip olduğunu göstermektedir.

\subsection{Öneri Motoru ve LLM Entegrasyonu}

Kural tabanlı motor ve LLM katmanı (Claude Sonnet + GPT-4o) birlikte 9 öneri üretmiştir: 2 KRİTİK ve 7 UYARI.
LLM'e iletilen zenginleştirilmiş bağlam; kapasite durumunu, log anomali özetini, SHAP özellik sıralamasını ve cascade risk skorunu kapsamaktadır.
Üretilen öneriler, tahmin modelinin sayısal çıktısını (tahmini eşik tarihi), SHAP açıklamasını (en etkili özellik) ve log analizini (erken uyarı korelasyonu) somut operasyonel aksiyonlarla birleştirmektedir.
Bu yapı, CloudMon'un reaktif bir izleme aracı olmanın ötesine geçerek açıklanabilir ve eyleme dönüştürülebilir öneriler üreten bir operasyonel zekâ sistemi olduğunu göstermektedir.


\section{Tezin Özgün Katkıları}
\label{sec:sonuc-katkilar}

Bu çalışma, kapsamlı deneysel değerlendirme çerçevesinde aşağıdaki dört özgün katkıyı sunmaktadır:

\begin{enumerate}
    \item \textbf{Log--metrik korelasyonuna dayalı kapasite tahmini:}
    Metrik ve log verilerinin ortak bir özellik uzayında birleştirilmesi, her iki veri kaynağının da kapasite tahminine anlamlı katkı yaptığını göstermiştir.
    SHAP analizi, log tabanlı özelliklerin yalnızca metrik özelliklerini değil, ağaç tabanlı modellerin gördüğü her iki veri tipini de geride bırakabildiğini nicel olarak ortaya koymuştur.
    Hibrit özellik uzayı, yalnızca metrik tabanlı yaklaşımlara kıyasla SHAP ilk 10'unda \%40 log özellik payıyla (vCPU 24\,s) üstün açıklanabilirlik sunmaktadır.

    \item \textbf{Dört aşamalı log işleme pipeline'ı:}
    Ham log sayımlarından ön-arıza örüntü eşleştirmesine ve cascade risk tespitine uzanan sistematik analiz hattı, tek bir aşamalı log sayımı yaklaşımlarına kıyasla daha zengin ve erken sinyaller üretmektedir.
    Ön-arıza örüntüsünün 68 saatlik öncü uyarı özelliği, operasyonel müdahale için pratik bir zaman penceresi sağlamaktadır.

    \item \textbf{Açıklanabilir kapasite tahmin sistemi:}
    Her tahmin için hesaplanan SHAP değerleri; hangi metriğin ya da log tipinin kararı yönlendirdiğini operatöre aktarmaktadır.
    Bu yaklaşım, kara kutu tahmin modellerini yorumlanabilir aksiyon önerilerine dönüştürerek bulut operatörlerinin karara güvenmesini ve gereksiz kota artışı kararlarını daha bilinçli değerlendirmesini sağlamaktadır.

    \item \textbf{Zenginleştirilmiş bağlamla LLM destekli öneri motoru:}
    Kapasite durumu, log anomali sinyalleri, SHAP özellik önem sıralaması ve cascade riski bir araya getirilerek yapılandırılmış bağlam oluşturulmuş; bu bağlam Claude Sonnet ve GPT-4o modellerine iletilmiştir.
    Üretilen öneriler, sayısal tahmin çıktısı ile yorumsal açıklamayı birleştiren servis bazında önceliklendirilmiş aksiyonlar içermekte; böylece operatörün ayrı sistemlere bakma zorunluluğunu ortadan kaldırmaktadır.
\end{enumerate}


\section{Sınırlılıklar}
\label{sec:sonuc-sinirliliklar}

Bu çalışmanın değerlendirilmesinde aşağıdaki sınırlılıklar göz önünde bulundurulmalıdır:

\begin{enumerate}
    \item \textbf{Eğitim kümesi büyüklüğü ve tek küme:}
    Modeller tek bir OpenStack kümesinin yaklaşık 9 aylık verisine (Temmuz 2025--Mart 2026) dayalı olarak eğitilmiştir.
    Bu durum, ağaç tabanlı modellerin zaman serisi genellemesinde başarısız olmasının kısmi nedeni olabilir; daha uzun eğitim veya farklı kümelerden elde edilen veriler ağaç modellerini güçlendirebilir.
    Ayrıca bulgular, yalnızca bu kümenin iş yükü örüntüsü için geçerli olup farklı ölçek veya kullanım biçimlerine sahip OpenStack kümelerine doğrudan genellenemeyebilir.

    \item \textbf{72 saatlik Disk tahmin gücünün sınırlı olması:}
    Disk 72\,s senaryosunda yalnızca MLP pozitif R\textsuperscript{2} (0{,}2212) üretmiş; diğer tüm modeller rassal tahminden daha kötü performans sergilemiştir.
    Disk kullanımının uzun vadede sergilediği düzensiz ve veri bağımlı örüntüsü, mevcut özellik kümesiyle yeterince yakalanamamaktadır.
    Bu tahmin ufku için ek özellik mühendisliği veya daha gelişmiş zaman serisi modelleri gerekmektedir.

    \item \textbf{LLM çıktılarının nesnel değerlendirmesi:}
    Öneri motorunun LLM katmanı, üretilen önerilerin kalitesini nesnel ölçütlerle değerlendirmek için bir insan değerlendirmesi (human evaluation) ya da karşılaştırma taban çizgisi (baseline) kullanılmamıştır.
    Mevcut değerlendirme, üretilen aksiyonların içerik kontrolüne dayanmakta olup operasyonel uygulamada etkinliklerini doğrulayan kontrollü bir deney yapılamamıştır.

    \item \textbf{Gerçek zamanlı pipeline yokluğu:}
    CloudMon, bu çalışma kapsamında toplu (batch) veri işleme ile test edilmiştir.
    Canlı OpenStack ortamında akış verisi işleme (streaming), model güncelleme gecikmesi ve düşük gecikimli öneri üretimi gibi operasyonel zorluklar deneysel ortamda ele alınmamıştır.
\end{enumerate}


\section{Gelecek Çalışma Önerileri}
\label{sec:sonuc-gelecek}

Bu çalışmanın bulguları ve sınırlılıkları ışığında aşağıdaki araştırma yönleri önerilmektedir:

\begin{enumerate}
    \item \textbf{Daha uzun ve çok kümeli veri ile model genellemesi:}
    Ağaç tabanlı modellerin genelleme yetersizliği, birden fazla OpenStack kümesinden toplanan ve birkaç yıllık zaman aralığını kapsayan verilerle yeniden değerlendirilebilir.
    Transfer öğrenme yaklaşımları, bir kümede eğitilen modelin diğer kümelere uyarlanmasını sağlayarak küme-özgü eğitim maliyetini azaltabilir.

    \item \textbf{Artımlı ve çevrimiçi öğrenme:}
    Bulut iş yükleri zamanla değiştiğinden, modellerin yeni gelen verilerle sürekli güncellenmesini sağlayan artımlı öğrenme (incremental learning) ve kavram kayması (concept drift) tespiti mekanizmaları sisteme entegre edilmelidir.
    Bu yaklaşım, modellerin uzun vadeli operasyonel geçerliliğini koruyacaktır.

    \item \textbf{Gerçek zamanlı akış mimarisi:}
    Mevcut toplu işleme pipeline'ı, Apache Kafka veya benzer bir akış platformu üzerine yeniden tasarlanarak düşük gecikmeli uyarı üretimi sağlanabilir.
    Gerçek zamanlı özellik hesaplama, özellikle cascade riski gibi zaman-kritik uyarıların pratikte değer üretmesini mümkün kılacaktır.

    \item \textbf{LLM çıktılarının kontrollü değerlendirmesi:}
    Öneri motorunun etkinliği, bulut operatörlerinin katılımıyla yürütülecek bir kullanıcı çalışmasıyla ölçülebilir.
    Bu çalışma; üretilen önerilerin doğruluğunu, uygulanabilirliğini ve operatör güvenini nicel olarak değerlendirme olanağı sunacaktır.
    Ayrıca farklı LLM modelleri (açık kaynak ile kapalı kaynak) karşılaştırmalı olarak test edilebilir.

    \item \textbf{Disk 72 saatlik tahmin için gelişmiş yöntemler:}
    Disk kullanımının uzun vadeli örüntüsünü daha iyi yakalamak amacıyla ayrıştırma tabanlı modeller (STL, Prophet) veya LSTM gibi dizi tabanlı derin öğrenme mimarileri araştırılabilir.
    Ayrıca I/O olaylarına dayalı olay güdümlü (event-driven) özellikler, bu senaryodaki tahmin boşluğunu kapatabilir.

    \item \textbf{Diğer bulut platformlarına genişleme:}
    CloudMon mimarisinin AWS, Azure veya Kubernetes kümelerine uyarlanması, platform bağımsız bir AIOps çerçevesi geliştirme ve farklı altyapı örüntülerinde bulguların tekrarlanabilirliğini test etme açısından değerli bir araştırma yönü oluşturmaktadır.
\end{enumerate}


\section{Genel Değerlendirme}

Bu tez, OpenStack tabanlı özel bulut altyapılarında proaktif kapasite yönetimini desteklemeye yönelik uçtan uca bir AIOps sistemi olan CloudMon'u tanıtmış ve deneysel olarak değerlendirmiştir.
Hibrit özellik mühendisliği, açıklanabilir makine öğrenmesi ve LLM tabanlı öneri üretimini entegre eden bu yaklaşım; mevcut literatürdeki izole çalışmaların ötesine geçerek birden fazla analiz katmanını tek bir operasyonel zekâ platformunda bir araya getirmektedir.

Deneysel bulgular, Doğrusal Regresyon'un 24 saatlik ufuklarda ağaç tabanlı modellere kıyasla belirgin biçimde üstün olduğunu; log özelliklerinin SHAP analizinde güçlü bir sinyal kaynağı oluşturduğunu ve dört aşamalı log pipeline'ının metrik eşik alarmlarına üç güne yakın zaman önünde uyarı üretme kapasitesine sahip olduğunu göstermiştir.
Bu bulgular bir arada değerlendirildiğinde, CloudMon'un bulut operasyonlarında önleyici müdahaleyi destekleyen pratik ve yorumlanabilir bir araç olma potansiyeli taşıdığı sonucuna ulaşılmaktadır.
